% Caption for output/publication/final/d4000_hminus_plane_with_jwst_data_epoch_matched.{pdf,png}
% Replace \citet{Lu2026} with the bibliography key of the Lu et al. (2026) sample paper.
\caption{%
As Fig.~\ref{fig:d4000_hminus_plane_with_jwst_data}, but with the mock population
restricted to the epochs that the JWST sample can occupy. The 19 quiescent galaxies of
\citet{Lu2026} lie at $z = 1.02$--1.96, where the age of a flat $\Lambda$CDM universe
($H_0 = 70$~km\,s$^{-1}$\,Mpc$^{-1}$, $\Omega_{\rm m} = 0.3$) is 3.3--5.7~Gyr; assuming
star formation begins at $t = 0$, only model epochs in this interval are kept (98\,000
of the 482\,000 epochs, and 3017 of the 5907 rapid-quenching epochs). Filled grey contours
enclose 68, 95 and 99.5 per cent of these epochs and the filled red contour 68 and 95 per
cent of their rapid-quenching class; dashed open contours show the corresponding
distributions over all epochs from 1 to 13~Gyr, as in
Fig.~\ref{fig:d4000_hminus_plane_with_jwst_data}. The blue curves trace the fiducial
history ($t_q = 3$~Gyr, $\tau_q = 0.3$~Gyr) at $\log Z/Z_\odot = -0.5$ to $+0.25$
within the same interval, with markers 0.5, 1 and 2~Gyr after quenching. No new models
are computed: the population is that of Fig.~\ref{fig:index_planes_agb_on_off}, with
D4000 measured at $R = 50$ and the H$^-$ bump at $R = 100$. Removing the old epochs
takes away the oldest, highest-D4000 part of both loci. Without TP-AGB light the
epoch-matched models reach at most $-0.017$~mag over the observed D4000 range, 0.02 to
0.06~mag weaker than every observed galaxy. With TP-AGB light 14 of the 19 galaxies lie
within the 99.5 per cent contour and three (one within the 68 per cent contour) within
the 95 per cent contour of the rapid-quenching class; the five outside are the four with
the weakest bumps ($-0.037$ to $-0.045$~mag at D4000 $= 1.6$--1.9), which fall below
the epoch-matched locus, and the highest-redshift, lowest-S/N object ($z = 1.96$,
D4000 $= 1.35$, $-0.073$~mag).%
}
